% Guion de exposición de PULSO (acompaña a la presentación v2).
% Compilar desde docs/latex/guion:
%   latexmk -lualatex guion.tex
\documentclass[11pt,a4paper,oneside]{article}

\usepackage{fontspec}
\usepackage[spanish,es-noshorthands,es-tabla,es-nodecimaldot]{babel}
\usepackage{amsmath,amssymb}
\setmainfont{TeX Gyre Pagella}
\setsansfont{TeX Gyre Heros}
\setmonofont{DejaVu Sans Mono}[Scale=0.82]
\usepackage[margin=2.3cm,headheight=14pt]{geometry}
\usepackage{graphicx}
\usepackage{booktabs,longtable,array}
\usepackage{xcolor}
\usepackage[most]{tcolorbox}
\usepackage{enumitem}
\usepackage{needspace}
\usepackage{titlesec}
\usepackage{fancyhdr}
\usepackage[hidelinks]{hyperref}

\graphicspath{{../figuras/}{../../../figures/}}

\definecolor{azul}{HTML}{2563A8}
\definecolor{rojo}{HTML}{C2362B}
\definecolor{verde}{HTML}{2F7D5B}
\definecolor{ambar}{HTML}{E0A100}
\definecolor{gris}{HTML}{5F6B77}

\titleformat{\section}{\sffamily\bfseries\Large\color{azul}}{}{0em}{}
\titleformat{\subsection}{\sffamily\bfseries\large}{}{0em}{}
\titlespacing*{\subsection}{0pt}{1.4em}{0.5em}

\pagestyle{fancy}
\fancyhf{}
\fancyhead[L]{\small\sffamily\color{gris}PULSO · guion de exposición}
\fancyhead[R]{\small\sffamily\color{gris}\nouppercase{\leftmark}}
\fancyfoot[C]{\small\thepage}
\renewcommand{\headrulewidth}{0.3pt}

\setlength{\parskip}{0.5em}
\setlength{\parindent}{0pt}
\setlist{itemsep=0.15em,topsep=0.25em}
\setlist[description]{font=\normalfont\bfseries,leftmargin=1.2em,labelindent=0em}

% Encabezado de cada diapositiva: número, título y tiempo sugerido.
\newcommand{\diapo}[3]{%
  \par\needspace{12\baselineskip}\vspace{0.8em}%
  \begin{tcolorbox}[colback=azul!8,colframe=azul,boxrule=0.6pt,arc=2pt,
    left=6pt,right=6pt,top=3pt,bottom=3pt,fontupper=\sffamily]
  \textbf{Diapositiva #1} \quad #2 \hfill {\small\color{gris}#3}
  \end{tcolorbox}\nopagebreak}
% Bloques del guion
\newtcolorbox{decir}{colback=white,colframe=gris!40,boxrule=0.4pt,arc=2pt,
  left=6pt,right=6pt,before upper={\setlength{\parskip}{0.55em}},title={\sffamily\small Qué decir},fonttitle=\color{black},
  coltitle=black,colbacktitle=gris!12,breakable}
\newtcolorbox{terminos}{colback=verde!5,colframe=verde,boxrule=0.5pt,arc=2pt,
  left=6pt,right=6pt,title={\sffamily\small Variables y términos},
  fonttitle=\color{black},colbacktitle=verde!15,breakable}
\newtcolorbox{preguntas}{colback=ambar!7,colframe=ambar!80!black,boxrule=0.5pt,
  arc=2pt,left=6pt,right=6pt,title={\sffamily\small Si preguntan},
  fonttitle=\color{black},colbacktitle=ambar!20,breakable}
\newcommand{\pr}[2]{\textbf{P.} \emph{#1}\par\textbf{R.} #2\par\smallskip}

\hypersetup{pdftitle={PULSO: guion de exposición},pdfauthor={Juan Diego Martínez Ayala}}

\begin{document}

\begin{titlepage}
\centering
\vspace*{3cm}
{\sffamily\bfseries\Huge\color{azul} PULSO\par}
\vspace{0.2cm}
{\sffamily\large\color{gris} Pagos, Umbrales, Liquidez y Sincronización Operativa\par}
\vspace{0.9cm}
{\sffamily\LARGE Guion de exposición\par}
\vspace{0.4cm}
{\large Texto de apoyo para cada diapositiva de la presentación (versión 2),\\
con la definición de cada variable, modelo, método y sigla\par}
\vfill
{\large Juan Diego Martínez Ayala\par}
{\normalsize Dirección: Jhon Hernán Díaz Forero\par}
{\normalsize Semillero de Modelamiento de Sistemas Complejos y Física Computacional\par}
\vspace{0.4cm}
{\normalsize Noviembre de 2026 · software v3.1.0\par}
\end{titlepage}

\tableofcontents
\clearpage

\section*{Cómo usar este guion}
\addcontentsline{toc}{section}{Cómo usar este guion}

La presentación (\texttt{pulso\_presentacion\_v2.pdf}) muestra pocos elementos por
diapositiva para respetar el tiempo de exposición. Este guion contiene lo que la
diapositiva no dice. Para cada diapositiva incluye:

\begin{description}
\item[Qué decir.] El texto de la exposición. Está escrito para decirse en voz alta y se
puede leer o parafrasear.
\item[Variables y términos.] La definición de cada símbolo de las ecuaciones y de cada
modelo, método o sigla que aparece en la diapositiva. Sirve para responder con
precisión.
\item[Si preguntan.] Preguntas probables del jurado o del público y una respuesta
breve y defendible.
\end{description}

El número de cada diapositiva es el que aparece en la esquina inferior derecha de la
presentación. Las páginas separadoras de sección no tienen número. Los tiempos
sugeridos suman unos 30 minutos.

\begin{center}
\small
\begin{tabular}{lll}
\toprule
Sección & Diapositivas & Tiempo \\
\midrule
Planteamiento & 1–5 & 5 min \\
Modelo unificado & 6 & 2 min \\
Tareas 1 a 4 & 7–10 & 5 min \\
Tarea 5 & 11–12 & 2.5 min \\
Tarea 6 & 13–17 & 6 min \\
Resultados & 18–25 & 7 min \\
Implementación & 26–28 & 1.5 min \\
Conclusiones & 29–31 & 1.5 min \\
\midrule
Total & & 30.5 min \\
\bottomrule
\end{tabular}
\end{center}

\clearpage

% ===================================================================
\section{Apertura y planteamiento}

\diapo{0}{Portada}{30 s}
\begin{decir}
Buenos días. Presento PULSO, sigla de Pagos, Umbrales, Liquidez y Sincronización
Operativa. Es el proyecto del Semillero de Modelamiento de Sistemas Complejos y Física
Computacional, dirigido por el profesor Jhon Hernán Díaz Forero.

El proyecto estudia cómo los problemas de liquidez se propagan en una red de pagos
digitales. También estudia si el ritmo con que llegan los pagos, y no solo su monto,
permite anticipar el impago de un crédito.
\end{decir}
\begin{terminos}
\begin{description}
\item[Liquidez.] Disponibilidad de dinero en el momento en que hay que pagar. Se
distingue de la \emph{solvencia}, que compara el total de activos con el total de
deudas. Una entidad solvente puede quedarse sin liquidez si el dinero le llega tarde.
\item[Umbral.] Nivel de déficit que una entidad puede absorber antes de dejar de pagar.
En el modelo representa el colchón de liquidez.
\item[Sincronización operativa.] Grado en que los ciclos de pago de las entidades
coinciden en el tiempo.
\end{description}
\end{terminos}

\diapo{1}{Contenido}{20 s}
\begin{decir}
La exposición sigue el orden en que se construyó el proyecto. Primero planteo el
problema y la pregunta. Luego presento el modelo y las cinco tareas de la propuesta del
semillero, que van de la red de pagos a la predicción con aprendizaje automático.
Después presento la tarea 6, que aplica todo al riesgo crediticio de un banco. Termino
con los resultados, la implementación y las conclusiones.
\end{decir}

\diapo{2}{Contexto}{1 min}
\begin{decir}
Hoy casi todos los pagos de una economía pasan por sistemas digitales: salarios,
facturas entre empresas, arriendos y cuotas de crédito. En Colombia el Banco de la
República puso en marcha Bre-B, un sistema de pagos inmediatos que conecta a todas las
entidades financieras.

Cada persona o empresa paga con el dinero que recibe de otras. Por eso los pagos forman
una red. En la red, cada nodo es un agente y cada flecha es un pago de un agente a otro.
La flecha tiene un peso, que es el monto, y una dirección, que va de quien paga a quien
recibe.

En el diagrama, A tiene un problema y deja de pagar. B y C pierden parte de su ingreso.
Si no tienen reservas suficientes, tampoco pueden pagar a D. El problema de A llega a D
aunque D no tenga ninguna relación directa con A.

La inmediatez es una ventaja para los usuarios, pero también hace que un problema de
liquidez se transmita más rápido. Este proyecto modela esa transmisión.
\end{decir}
\begin{terminos}
\begin{description}
\item[Sistema de pagos.] Infraestructura que transfiere dinero entre cuentas de
distintas entidades y liquida las obligaciones.
\item[Bre-B.] Sistema de pagos inmediatos de bajo valor del Banco de la República de
Colombia, interoperable entre bancos y billeteras. En el proyecto aparece solo como
escenario de referencia. Los parámetros no están calibrados con sus datos.
\item[Red dirigida y ponderada.] Conjunto de nodos unidos por flechas con dirección y
peso. Se representa con una matriz $W$, donde $W_{ij}$ es el monto que $j$ paga a $i$.
\item[Impago.] Incumplimiento de una obligación de pago.
\end{description}
\end{terminos}
\begin{preguntas}
\pr{¿El proyecto usa datos de Bre-B?}{No. Todos los datos son simulados. Bre-B sirve
como motivación y como escenario ilustrativo de la aplicación (una red núcleo-periferia
de 245 entidades), con parámetros que no están calibrados.}
\end{preguntas}

\diapo{3}{Ciclos de pago y desfase}{1 min}
\begin{decir}
Los pagos no llegan de manera continua: siguen ciclos. La nómina llega cada quince o
treinta días, las facturas tienen plazos y las cuotas tienen fechas fijas.

Lo importante es la coordinación. Si el salario llega el día 30 y la cuota vence el día
5, el cliente tiene el dinero a tiempo. Si la nómina se corre al día 10, el monto es el
mismo pero el cliente necesita un colchón para cubrir esos días. A ese desfase lo llamamos
desincronización.

La gráfica de la izquierda muestra el mismo cliente con el mismo ruido, con y sin costo
por desfase. Con desfase, el déficit es mayor en todo momento. La gráfica de la derecha
muestra un empleado cuyo ritmo natural es 8~\% más rápido que el de su nómina. Con un
acoplamiento fuerte, el empleado ajusta su ritmo y se sincroniza; con un acoplamiento
débil, se desfasa cada vez más.

La idea central del proyecto es que el ingreso desfasado genera estrés financiero aunque
el monto no cambie.
\end{decir}
\begin{terminos}
\begin{description}
\item[Ciclo de pagos.] Periodo que se repite entre dos pagos del mismo tipo, por ejemplo
un mes de nómina.
\item[Fase.] Posición dentro del ciclo, medida como un ángulo entre $0$ y $2\pi$. La
fase $0$ es el inicio del ciclo y $2\pi$ es su final.
\item[Desfase.] Diferencia entre la fase de quien paga y la de quien recibe.
\item[Déficit.] Faltante de liquidez en un momento dado. En el modelo es la variable
$x_i$.
\item[Colchón.] Reserva que permite cubrir un déficit temporal sin dejar de pagar.
\item[Acoplamiento.] Intensidad con que un agente ajusta su ritmo al de quienes le
pagan.
\end{description}
\end{terminos}

\diapo{4}{Limitación del score clásico y pregunta de investigación}{1 min}
\begin{decir}
Los bancos deciden a quién prestar con un score crediticio. El score es un modelo
estadístico que estima la probabilidad de impago con variables del propio cliente:
ingreso, endeudamiento, historial de mora, antigüedad y tipo de empleo.

El score evalúa a cada cliente de forma aislada. No ve tres cosas: si quienes le pagan
al cliente están en problemas, si todo su ingreso depende de un solo pagador, y si sus
ingresos llegan desfasados de sus obligaciones.

Esa información existe en los registros de transacciones del banco, pero el score
clásico no la usa.

La pregunta del proyecto es esta: la sincronización de los ciclos de pago, estimada
únicamente con las transacciones, ¿aporta información sobre el impago que no esté ya en
el score clásico ni en la estructura estática de la red? La comparación con la red
estática es la parte exigente: ya existen trabajos que usan la red de un cliente para
mejorar el score. Lo nuevo es la dinámica.
\end{decir}
\begin{terminos}
\begin{description}
\item[Score crediticio.] Modelo estadístico que asigna a cada cliente una probabilidad
de impago, normalmente con regresión logística sobre variables individuales.
\item[PD.] Probabilidad de incumplimiento (\emph{probability of default}): probabilidad
de que el cliente deje de pagar su crédito en un horizonte, aquí doce meses.
\item[Red estática.] Variables que describen quién le paga a quién y cuánto, sin
considerar el momento de los pagos: concentración del ingreso, número de pagadores,
centralidad.
\item[Red dinámica.] Variables que describen el ritmo: sincronización, desfase,
regularidad de los ciclos.
\end{description}
\end{terminos}
\begin{preguntas}
\pr{¿No es esto lo que ya hacen los bancos con el historial crediticio?}{El historial
describe al cliente por sí solo. Las variables de red estática describen sus relaciones
y ya aparecen en la literatura (por ejemplo, Óskarsdóttir y Bravo, 2021). La
sincronización de los ciclos de pago no aparece en esos trabajos, y por eso el
experimento la compara contra ambos niveles.}
\end{preguntas}

\diapo{5}{Hipótesis y objetivos}{1 min}
\begin{decir}
La hipótesis es falsable, es decir, puede resultar falsa. Si el ingreso desfasado
genera déficit, las variables de sincronización deben mejorar la predicción del impago.
Si ese mecanismo no existe, no deben mejorarla. Por eso el experimento incluye un
control nulo: un mundo idéntico donde la red no genera déficit. Si las variables
mejoraran la predicción también ahí, el resultado sería un artefacto del método.

El objetivo general es construir, paso a paso, un modelo físico-matemático de la red de
pagos y aplicarlo al riesgo crediticio. La tabla muestra las seis tareas. Las cinco
primeras vienen de la propuesta del semillero: déficit en la red, difusión,
sincronización, impagos en cascada y aprendizaje automático. La sexta es la aplicación a
la cartera de un banco.
\end{decir}
\begin{terminos}
\begin{description}
\item[Hipótesis falsable.] Afirmación que predice un resultado observable que podría no
ocurrir; si no ocurre, la hipótesis se rechaza.
\item[Control nulo.] Experimento idéntico al principal en el que se apaga el mecanismo
estudiado. Mide cuánto ``efecto'' produce el método por sí solo.
\item[Kuramoto.] Modelo de osciladores acoplados propuesto por Yoshiki Kuramoto (1975)
para estudiar la sincronización. Se explica en la diapositiva 9.
\end{description}
\end{terminos}

% ===================================================================
\section{Modelo unificado}

\diapo{6}{Un modelo, seis tareas}{2 min}
\begin{decir}
Esta es la ecuación central del proyecto. Cada entidad $i$ tiene dos variables: su
déficit $x_i$ y la fase de su ciclo de pagos $\theta_i$.

La primera ecuación dice cómo cambia el déficit en el tiempo. Tiene cinco términos.
Primero, $-\alpha x_i$: cada entidad corrige su déficit por sí misma, por ejemplo con
ahorro o crédito de corto plazo. Segundo, $\beta\sum_j W_{ij}x_j$: si quienes le pagan
están en déficit, le pagan menos y su déficit aumenta; este es el contagio. Tercero,
$-D_0(Lx)_i$: la redistribución de liquidez entre entidades, que traslada déficit sin
crearlo. Cuarto, $\gamma k_i(1-r_i)$: el ingreso que llega desfasado se vuelve déficit.
Quinto, $\eta_i(t)$: los choques externos.

La segunda ecuación describe el ritmo de pagos. Cada entidad tiene su ritmo natural
$\omega_i$ y lo ajusta según el ritmo de quienes le pagan, con intensidad $K$.

Finalmente, si el déficit supera el umbral $x_i^{(c)}$, la entidad entra en impago y deja
de pagar.

Las etiquetas T1 a T4 indican en qué tarea entra cada término. El criterio de diseño es
que, si se apagan los términos nuevos, el modelo coincide exactamente con el de la tarea
anterior. Eso se verificó con pruebas automáticas. Así, cada resultado se puede
comparar con el anterior.
\end{decir}
\begin{terminos}
\begin{description}
\item[$x_i(t)$.] Déficit de liquidez de la entidad $i$ en el instante $t$. En las
tareas 1–5 se mide en ``pagos medios por ciclo''.
\item[$t$.] Tiempo. En las tareas 1–5, una unidad es un ciclo de pagos (un día hábil);
en la tarea 6, un mes.
\item[$dx_i/dt$.] Tasa de cambio del déficit: cuánto crece o decrece por unidad de
tiempo.
\item[$\alpha$.] Tasa de autocorrección [1/ciclo]. Con solo este término el déficit
decae como $e^{-\alpha t}$; $1/\alpha$ es el tiempo típico de corrección.
\item[$\beta$.] Intensidad del contagio [1/(ciclo·pago)]: déficit que genera en $i$
cada unidad de déficit de un pagador, por unidad de pago.
\item[$W_{ij}$.] Pago medio de $j$ hacia $i$ por ciclo. Las filas de $W$ son ingresos
y las columnas son egresos.
\item[$\sum_j W_{ij}x_j$.] Suma del déficit de todos los pagadores de $i$, ponderada
por cuánto le paga cada uno.
\item[$D_0$.] Coeficiente de redistribución [1/(ciclo·pago)].
\item[$L$.] Laplaciano de flujo, $L=D_\text{out}-W$ (diapositiva 8).
\item[$(Lx)_i$.] Componente $i$ del vector $Lx$.
\item[$\gamma$.] Peso del desfase en el déficit [1/ciclo].
\item[$k_i$.] Ingreso por la red: $k_i=\sum_jW_{ij}$.
\item[$r_i$.] Coherencia local: entre $-1$ (ingreso en contrafase) y $1$ (ingreso en
fase).
\item[$\eta_i(t)$.] Forzamiento externo: choques de liquidez en las tareas 1–5 y ruido
aleatorio en la tarea 6.
\item[$\theta_i$.] Fase del ciclo de pagos de $i$, en radianes.
\item[$\omega_i$.] Frecuencia natural [rad/ciclo]: ritmo de $i$ si no tuviera pagadores.
\item[$K$.] Intensidad del acoplamiento de fase [rad/ciclo].
\item[$\langle k\rangle$.] Ingreso medio por entidad. Normaliza $K$ (diapositiva 9).
\item[$\sin(\theta_j-\theta_i)$.] Positivo si el pagador $j$ va adelantado: acelera a
$i$. Negativo si va atrasado: lo frena.
\item[$x_i^{(c)}$.] Umbral de impago (colchón).
\item[Criterio de anidamiento.] Con $D_0=\gamma=K=0$ y sin umbral se recupera la
tarea 1; con $\gamma=K=0$, la tarea 2; y así sucesivamente.
\end{description}
\end{terminos}
\begin{preguntas}
\pr{¿Por qué un modelo lineal en $x$?}{Es la aproximación más simple que conserva el
contagio y permite resultados analíticos: el umbral de estabilidad $R=1$, la
conservación bajo difusión y la tasa de recuperación. La no linealidad entra por el
umbral y por las fases. Es una limitación declarada.}
\pr{¿El déficit puede afectar el ritmo de pagos?}{En esta versión no: el acople va del
ritmo al déficit. La retroalimentación inversa aparece en el trabajo futuro.}
\end{preguntas}

% ===================================================================
\section{Tareas 1 a 4: de la red a las cascadas}

\diapo{7}{Tarea 1: red de pagos e índice de amplificación}{1.5 min}
\begin{decir}
La primera tarea construye la red y estudia solo dos términos: la autocorrección y el
contagio.

Generamos redes con tres formas. En la de Erdős–Rényi, cada par de entidades se conecta
al azar con la misma probabilidad; sirve de referencia. En la de Barabási–Albert, las
entidades nuevas prefieren conectarse con las que ya tienen muchas conexiones; aparecen
unos pocos nodos muy conectados. En la núcleo-periferia, un grupo pequeño de grandes
entidades está muy conectado entre sí y las demás se conectan con ese núcleo; es la
forma que se observa en los sistemas de pagos reales. Los montos siguen una distribución
lognormal: muchos pagos pequeños y pocos muy grandes.

Con solo estos dos términos, la ecuación es lineal: $\dot x=(-\alpha I+\beta W)x$. Como
todos los pagos son positivos, el teorema de Perron–Frobenius garantiza que el valor
propio dominante de $W$ es real y positivo; lo llamamos radio espectral $\rho(W)$. El
sistema es estable, es decir, el déficit se extingue después de un choque, si y solo si
$R=\beta\rho(W)/\alpha$ es menor que 1.

$R$ funciona como el número reproductivo de una epidemia: compara cuánto déficit genera
la red con cuánto corrige cada entidad. Si $R<1$, la red amortigua. Si $R>1$, la red
amplifica cualquier choque.

La figura muestra la tasa de crecimiento del déficit en el plano $\alpha$–$\beta$ para
las tres topologías. Azul es estable, rojo inestable. La línea negra es la frontera
teórica $\beta=\alpha/\rho(W)$ y coincide con el cambio de color. Las redes más
concentradas tienen $\rho$ mayor y pierden la estabilidad con menos contagio.

Además, los vectores propios dominantes ordenan a las entidades: el izquierdo mide
cuánto déficit propaga cada una (importancia) y el derecho cuánto acumula
(vulnerabilidad).
\end{decir}
\begin{terminos}
\begin{description}
\item[Erdős–Rényi (ER).] Red aleatoria en la que cada arista existe con probabilidad
$p$, independiente de las demás. El número de conexiones por nodo es similar para
todos.
\item[Barabási–Albert (BA).] Red que crece por conexión preferencial: un nodo nuevo se
une a los existentes con probabilidad proporcional a su número de conexiones. Produce
nodos concentradores (\emph{hubs}).
\item[Núcleo-periferia.] Red con un núcleo denso y una periferia que se conecta casi
solo al núcleo.
\item[Distribución lognormal.] Distribución de una variable cuyo logaritmo es normal.
Montos $W_{ij}=\sqrt{s_is_j}\,\varepsilon_{ij}$, con tamaños $s$ y ruido $\varepsilon$
lognormales; luego se normaliza a peso medio 1.
\item[$I$.] Matriz identidad.
\item[Valor propio y vector propio.] $\lambda$ y $v\neq0$ tales que $Wv=\lambda v$. Los
valores propios de $-\alpha I+\beta W$ son $-\alpha+\beta\lambda_k(W)$.
\item[Teorema de Perron–Frobenius.] Una matriz con entradas no negativas tiene un valor
propio real igual a su radio espectral, con un vector propio de entradas no negativas.
\item[$\rho(W)$.] Radio espectral: el mayor módulo de los valores propios de $W$.
\item[$R$.] Índice de amplificación $\beta\rho(W)/\alpha$. Adimensional.
\item[Número reproductivo ($R_0$).] En epidemiología, número medio de contagios que
produce un infectado. La analogía es exacta en forma: umbral en 1.
\item[Importancia y vulnerabilidad.] Componentes de los vectores propios izquierdo
($u^\top W=\rho u^\top$) y derecho ($Wv=\rho v$).
\end{description}
\end{terminos}
\begin{preguntas}
\pr{¿Cómo se verificó?}{Se comparó la integración numérica con la solución exacta
$x(t)=e^{Mt}x_0$; el error relativo es menor que $10^{-6}$. Con $R=0.9$ el déficit se
extingue y con $R=1.1$ crece, como predice la teoría.}
\end{preguntas}

\diapo{8}{Tarea 2: redistribución y estabilidad}{1.5 min}
\begin{decir}
La segunda tarea agrega la redistribución de liquidez. Usamos el Laplaciano de flujo
$L=D_\text{out}-W$, donde $D_\text{out}$ es una matriz diagonal con lo que paga cada
entidad. El término $-D_0Lx$ traslada déficit de una entidad a otra a lo largo de los
pagos.

Las columnas de $L$ suman cero. Eso implica que la suma de todos los déficits no cambia
por este término: el déficit se mueve, pero no se crea ni se destruye. Lo verificamos
numéricamente: la suma cambia menos de $10^{-10}$.

Con los tres términos, la matriz del sistema es $M=-\alpha I-D_0L+\beta W$. Su valor
propio con mayor parte real, $\sigma$, decide la estabilidad: si $\sigma<0$ la red se
recupera, y el tiempo de recuperación es $\tau_\text{rec}=-1/\sigma$.

La figura de la izquierda muestra que las redes con mayor conectividad algebraica
$\lambda_2$ se recuperan más rápido. En redes simétricas la tasa de recuperación es
exactamente $\alpha+D_0\lambda_2$, y la simulación coincide con un error del 1~\%.

La figura de la derecha muestra comunidades: grupos de entidades que se pagan mucho
entre sí y poco con el resto. Las obtenemos con el vector de Fiedler del Laplaciano
normalizado.
\end{decir}
\begin{terminos}
\begin{description}
\item[$D_\text{out}$.] Matriz diagonal con $\mathrm{out}_j=\sum_iW_{ij}$, el pago total
de $j$ por ciclo.
\item[Laplaciano de flujo $L$.] $L=D_\text{out}-W$. Como $\sum_iL_{ij}=\mathrm{out}_j-
\sum_iW_{ij}=0$, se cumple $\mathbf 1^\top L=0$ y $\frac{d}{dt}\sum_ix_i=0$ bajo este
término.
\item[$D_0$.] Velocidad de la redistribución.
\item[$M$.] Matriz del sistema lineal completo.
\item[$\sigma$.] Máxima parte real de los valores propios de $M$. Negativa: estable.
\item[$\tau_\text{rec}$.] Tiempo de recuperación $-1/\sigma$, en ciclos: tiempo en que
una perturbación se reduce en un factor $e$.
\item[$\lambda_2$ (conectividad algebraica).] Segundo valor propio más pequeño del
Laplaciano de la red simetrizada $S=(W+W^\top)/2$. Vale 0 si la red está partida en
dos y crece con la conectividad.
\item[Vector de Fiedler.] Vector propio asociado a $\lambda_2$. El signo de sus
componentes divide la red en dos grupos poco conectados.
\item[Laplaciano normalizado.] $I-D^{-1/2}SD^{-1/2}$, con $D$ la matriz de grados de
$S$. Implementa el corte normalizado de Shi y Malik (2000), que evita separar solo
nodos periféricos.
\item[Bisección recursiva.] Se divide en dos la comunidad más grande hasta obtener el
número pedido de comunidades.
\end{description}
\end{terminos}
\begin{preguntas}
\pr{¿Qué representa la redistribución en la práctica?}{Mecanismos que trasladan presión
de liquidez sin crearla: líneas de crédito interbancario, compensación de pagos. Es una
idealización.}
\end{preguntas}

\diapo{9}{Tarea 3: sincronización de los ciclos de pago}{1.5 min}
\begin{decir}
La tercera tarea modela el ritmo de los pagos con el modelo de Kuramoto, el modelo
estándar de sincronización en física. Cada entidad es un oscilador con una fase
$\theta_i$ y un ritmo natural $\omega_i$. Si el acoplamiento $K$ es grande, las fases se
alinean, como cuando un público que aplaude termina aplaudiendo al unísono.

Medimos la sincronización con el parámetro de orden $r$, que vale 0 si las fases están
dispersas y 1 si todas coinciden. Para cada entidad definimos la coherencia local
$r_i$: el promedio del coseno del desfase con sus pagadores, ponderado por lo que le paga
cada uno.

Este término se conecta con el déficit mediante $\gamma k_i(1-r_i)$. Si el ingreso
llega en fase, $r_i=1$ y el término vale cero. Si llega en contrafase, $r_i=-1$ y el
término es máximo. Así la sincronización deja de ser descriptiva y afecta la liquidez.

Para validar, usamos un resultado exacto: en una red completa con frecuencias de
Lorentz, $r=\sqrt{1-2\Delta/K}$. Con $K=2$ el valor teórico es 0.7071 y el simulado
0.7076. La figura muestra que las redes dispersas se sincronizan de forma más gradual
que la red completa.

Una decisión técnica: normalizamos el acoplamiento por el ingreso medio y no por el
número de entidades. En redes dispersas, dividir por $N$ anula el acoplamiento cuando la
red crece.
\end{decir}
\begin{terminos}
\begin{description}
\item[Modelo de Kuramoto.] $\dot\theta_i=\omega_i+\frac{K}{\langle k\rangle}\sum_jW_{ij}
\sin(\theta_j-\theta_i)$. Cada oscilador acelera o frena según la diferencia de fase con
quienes lo influyen.
\item[$\omega_i$ y distribución de Lorentz.] Frecuencias naturales sorteadas de una
distribución de Lorentz (o de Cauchy) con centro $\bar\omega=2\pi$ (un ciclo por día) y
semiancho $\Delta$.
\item[$\Delta$.] Dispersión de los ritmos naturales. Mayor $\Delta$, más difícil
sincronizar.
\item[$r$.] Parámetro de orden global $r=\bigl|\frac1N\sum_je^{i\theta_j}\bigr|\in[0,1]$.
Aquí $i$ dentro del exponente es la unidad imaginaria.
\item[$r_i$.] Coherencia local $\frac1{k_i}\sum_jW_{ij}\cos(\theta_j-\theta_i)\in[-1,1]$.
\item[$\gamma k_i(1-r_i)$.] Ingreso desfasado convertido en déficit; es igual a
$\gamma\sum_jW_{ij}(1-\cos(\theta_j-\theta_i))$. Unidades de dinero por ciclo, como
los demás términos.
\item[$K_c=2\Delta$.] Acoplamiento crítico en la red completa con Lorentz: por debajo,
$r\approx0$.
\item[Normalización por $\langle k\rangle$.] Rodrigues et al.\ (2016). En la red
completa coincide con el $K/N$ clásico.
\end{description}
\end{terminos}
\begin{preguntas}
\pr{¿Por qué Kuramoto y no otro modelo?}{Es el modelo mínimo con una transición de
sincronización conocida analíticamente. Permite validar el código contra un resultado
exacto y tiene una interpretación directa: cada entidad ajusta su ciclo al de quienes le
pagan.}
\end{preguntas}

\diapo{10}{Tarea 4: impagos y cascadas}{1.5 min}
\begin{decir}
La cuarta tarea introduce el impago. Cada entidad tiene un colchón proporcional a lo que
paga, con algo de variación aleatoria. Si su déficit supera el colchón, la entidad deja
de pagar. Sus receptores pierden ese ingreso, su déficit salta, y algunos pueden cruzar
su propio umbral. Repetimos el proceso en el mismo instante hasta que no hay impagos
nuevos: eso es una cascada.

Simulamos dos choques. El retiro de liquidez es un salto en el déficit de una entidad.
La caída operativa es una entidad que recibe pero no paga durante un tiempo.

Comparamos con tres modelos clásicos de contagio financiero: Eisenberg–Noe, DebtRank y
Watts.

La figura muestra la ventana de contagio. Con pocas conexiones, el choque no tiene por
dónde propagarse. Con muchas conexiones, cada entidad recibe de muchos pagadores y la
pérdida de uno se diluye. Las cascadas son máximas con conectividad intermedia. Nuestra
dinámica produce esa ventana y difiere de DebtRank y de Watts, que no tienen dinámica
temporal ni recuperación.
\end{decir}
\begin{terminos}
\begin{description}
\item[$x_i^{(c)}=c\,\mathrm{out}_i\,\varepsilon_i$.] Umbral: $c$ es el colchón medido en
ciclos de pagos salientes, $\mathrm{out}_i$ lo que paga $i$ por ciclo y $\varepsilon_i$
un factor lognormal de media 1.
\item[$W_{\cdot i}\leftarrow0$.] La columna $i$ de $W$ se hace cero: $i$ deja de pagar
a todos.
\item[$x_k\mathrel{+}=\lambda\,\mathrm{out}_k W_{ki}/k_k$.] Salto en el déficit de cada
receptor $k$: fracción de su ingreso que perdió ($W_{ki}/k_k$) por su escala de pagos
($\mathrm{out}_k$) por $\lambda$, el parámetro de pérdida.
\item[Cascada instantánea.] Repetición de la regla de impago en el mismo instante hasta
que no hay impagos nuevos.
\item[Orden en cada paso.] Eventos programados, paso de integración, detección de
impagos y cascada.
\item[Eisenberg–Noe (2001).] Modelo de liquidación: cada entidad paga el mínimo entre lo
que debe y lo que tiene, $p=\min(\bar p,e+\Pi^\top p)$. $\bar p$ son las obligaciones,
$e$ los activos externos y $\Pi$ la matriz de proporciones de deuda.
\item[DebtRank.] Battiston et al.\ (2012). Propaga estrés de acreedor a deudor y cada
nodo transmite una sola vez, sin reverberación.
\item[Modelo de Watts (2002).] Un nodo falla cuando la fracción de vecinos caídos supera
un umbral $\phi$.
\item[Colapso global.] Impago de al menos el 30~\% de las entidades.
\end{description}
\end{terminos}

% ===================================================================
\section{Tarea 5: predicción sobre simulaciones}

\diapo{11}{Diseño del experimento numérico}{1 min}
\begin{decir}
Una simulación es precisa pero lenta para responder en tiempo real. La quinta tarea
entrena modelos de aprendizaje automático sobre miles de simulaciones para predecir el
resultado en milisegundos.

Generamos 300 redes y 20\,000 escenarios. Cada escenario combina una red con valores de
13 parámetros, entre ellos los de la tabla, el tamaño y el tipo de choque. Los valores se
eligen con un hipercubo latino, que cubre el espacio de parámetros de manera uniforme.
Además, concentramos más escenarios cerca de $R=1$, donde el comportamiento cambia y
predecir es más difícil.

Construimos dos tablas: una por escenario, con la etiqueta de colapso global, y una por
entidad, con el tamaño de la cascada que produce un choque en esa entidad.

Para evaluar, usamos GroupKFold por red: ninguna red aparece a la vez en entrenamiento y
en prueba. Si no se hiciera así, el modelo podría memorizar la red y el resultado sería
optimista.
\end{decir}
\begin{terminos}
\begin{description}
\item[Hipercubo latino.] Muestreo en el que el rango de cada parámetro se divide en $N$
intervalos y cada intervalo se usa exactamente una vez. Cubre mejor el espacio que un
muestreo aleatorio simple.
\item[Muestreo cerca de $R=1$.] $R=1+0.6\,\mathrm{sign}(u)|u|^2$, con $u$ uniforme en
$[-1,1]$.
\item[Los 13 parámetros.] $\alpha$, $R$, $D_0$, $K$, $\gamma$, $\Delta$, colchón $c$,
pérdida $\lambda$, tamaño del choque, duración, tipo de choque, entidad chocada y red.
\item[GroupKFold.] Validación cruzada en $k$ particiones (aquí 5) en la que todas las
filas de un mismo grupo (una red) quedan en la misma partición.
\item[Validación cruzada.] Se entrena con $k-1$ particiones y se evalúa en la restante,
rotando; cada predicción se hace sobre datos que el modelo no vio.
\item[Semilla por fila.] Cada escenario tiene su propia semilla derivada con
\texttt{SeedSequence}: el conjunto de datos se regenera idéntico con cualquier número de
procesos.
\end{description}
\end{terminos}

\diapo{12}{Tarea 5: resultados}{1.5 min}
\begin{decir}
Probamos cuatro modelos. La regresión logística es la referencia interpretable. XGBoost
combina cientos de árboles de decisión. La GCN es una red neuronal que trabaja sobre el
grafo. La PINN es una red neuronal que respeta las ecuaciones del modelo.

Para el colapso global, XGBoost alcanza un AUC de 0.98. Para el tamaño de la cascada
por entidad, la GCN explica el 60~\% de la variación. La figura muestra el AUC del
colapso por tramos de $R$ y por topología. XGBoost se mantiene entre 0.97 y 0.99 en todos
los grupos. La regresión logística baja a 0.92 en la zona más inestable, $R>1.2$, porque
allí la relación entre las variables y el colapso deja de ser lineal.

Una simulación tarda 0.08 segundos; una predicción, 0.2 milisegundos. Eso permite el
mapa de riesgo instantáneo de la aplicación.

La PINN resuelve el problema inverso: a partir de trayectorias observadas, estima los
parámetros. Recupera $\gamma$ y $K$ con error menor al 1~\%. El parámetro $\beta$ tiene
más error, 23~\%, porque su efecto se confunde con el de $\alpha$.

Una limitación: dentro de una misma red, el orden de las entidades por riesgo se predice
solo de forma moderada, con una correlación de Spearman de 0.3. Por eso la aplicación
permite verificar cualquier predicción con una simulación.
\end{decir}
\begin{terminos}
\begin{description}
\item[Regresión logística.] Modelo $P(y=1)=1/(1+e^{-(b_0+\sum_kb_kz_k)})$. Cada
coeficiente indica cuánto cambia el logaritmo de la razón de probabilidades por unidad
de la variable.
\item[XGBoost.] \emph{Extreme Gradient Boosting}. Suma de árboles de decisión entrenados
en secuencia; cada árbol corrige los errores de los anteriores (\emph{gradient
boosting}).
\item[Árbol de decisión.] Modelo que divide los datos con preguntas del tipo ``¿$R>1$?''
hasta llegar a una predicción.
\item[GCN.] \emph{Graph Convolutional Network} (Kipf y Welling, 2017). Cada capa promedia
las variables de los vecinos de cada nodo y aplica una transformación aprendida. Se
usan 3 capas con una conexión residual que lleva las variables propias del nodo hasta
la salida.
\item[PINN.] \emph{Physics-Informed Neural Network} (Raissi et al., 2019). Red neuronal
entrenada para ajustar los datos y, a la vez, cumplir las ecuaciones del modelo. Los
parámetros físicos son incógnitas del entrenamiento.
\item[Problema inverso.] Estimar parámetros a partir de observaciones de la solución.
\item[AUC.] Área bajo la curva ROC. Probabilidad de que el modelo asigne más riesgo a un
caso positivo elegido al azar que a uno negativo. 0.5 es azar; 1 es perfecto.
\item[Curva ROC.] Tasa de verdaderos positivos contra tasa de falsos positivos para todos
los umbrales de decisión.
\item[PR-AUC.] Área bajo la curva precisión-exhaustividad. Más informativa cuando los
positivos son pocos.
\item[$R^2$.] Coeficiente de determinación: fracción de la varianza de la variable
objetivo que explica el modelo.
\item[Spearman ($\rho_s$).] Correlación entre los órdenes (rangos) de dos variables.
\end{description}
\end{terminos}
\begin{preguntas}
\pr{¿Para qué predecir si se puede simular?}{Para explorar en tiempo real: la predicción
es unas 400 veces más rápida por escenario y unas 20\,000 veces más rápida para la red
completa de entidades. La simulación queda como verificación.}
\pr{¿Por qué $\beta$ tiene más error en la PINN?}{En la parte lineal, $\alpha$ y $\beta$
entran combinadas en $M$: varias parejas producen trayectorias casi iguales. Es un
problema de identificabilidad, no del código.}
\end{preguntas}

% ===================================================================
\section{Tarea 6: riesgo crediticio multicapa}

\diapo{13}{De entidades a clientes de un banco}{1 min}
\begin{decir}
Hasta aquí las entidades eran abstractas. La tarea 6 lleva el modelo a la cartera de un
banco: los nodos son personas y empresas clientes del banco.

No existen datos públicos con la red de pagos de los clientes de un banco. Por eso
generamos una cartera sintética con un mecanismo conocido y verificamos que el método lo
encuentra. Es el mismo principio con que se valida un instrumento de medida: primero se
mide algo cuyo valor se conoce.

La cartera se construye para producir exactamente lo que un banco tendría: los atributos
de cada cliente y el registro de sus transacciones. El banco no ve el modelo interno; solo
ve los pagos.

Las relaciones están en tres capas. Nómina: la empresa paga a su empleado cada quince o
treinta días; es la dependencia más fuerte. Proveedores: una empresa paga a otra según
sus facturas. Transferencias: pagos entre personas, irregulares y de menor monto.
\end{decir}
\begin{terminos}
\begin{description}
\item[Cartera.] Conjunto de clientes con crédito en un banco.
\item[Datos sintéticos.] Datos generados por simulación con reglas conocidas.
\item[Red multicapa.] Red con varios tipos de relación entre los mismos nodos; cada tipo
es una capa con su propia matriz $W^\ell$ (Boccaletti et al., 2014).
\item[F y P.] Firma (empresa) y persona.
\item[Capa N, S, T.] Nómina, proveedores (\emph{suppliers}) y transferencias. En el
registro aparece además la capa C: cuotas de crédito pagadas al banco.
\end{description}
\end{terminos}
\begin{preguntas}
\pr{¿Qué valor tiene un resultado con datos sintéticos?}{Demuestra que, si el mecanismo
existe, el método lo detecta con la información que un banco tiene, y que no detecta
nada cuando el mecanismo no existe. No demuestra que el mecanismo exista en una cartera
real; eso requiere datos reales y es el siguiente paso.}
\pr{¿Hay trabajos similares?}{Letizia y Lillo (2019) usan redes de pagos entre empresas
para calificación crediticia; Tobback et al.\ (2017) y Óskarsdóttir y Bravo (2021) usan
redes relacionales y multicapa para predecir impago. Ninguno usa la sincronización de los
ciclos de pago.}
\end{preguntas}

\diapo{14}{Cartera sintética}{1.5 min}
\begin{decir}
La cartera de referencia tiene 2000 personas y 100 empresas. El 55~\% de las personas
son informales: no reciben nómina de una empresa del banco. La cifra es del orden de la
informalidad laboral reportada por el DANE.

El tamaño de las empresas sigue una distribución de Pareto: pocas empresas grandes y
muchas pequeñas. Cada empresa paga nómina mensual o quincenal. Los empleados formales
heredan el ritmo de su empleador con una variación del 8~\%; los informales tienen
ritmos más dispersos, con una variación del 25~\%.

La figura es un ejemplo reducido para ver las capas: los cuadrados son empresas y los
puntos son personas.

Cada cliente tiene una probabilidad de impago verdadera. Se calcula con una función
logística de sus variables clásicas más un componente oculto que el banco no ve. Ese
componente se revela en parte a través del historial de mora, como en la realidad. Las
columnas con la verdad del modelo empiezan por \texttt{truth\_} y el código impide que
cualquier modelo las use: si alguien lo intenta, el programa se detiene con un error.
\end{decir}
\begin{terminos}
\begin{description}
\item[Informalidad.] Trabajo sin vínculo laboral formal. En el modelo: persona sin
empleador dentro de la red del banco.
\item[DANE.] Departamento Administrativo Nacional de Estadística de Colombia.
\item[Distribución de Pareto.] Distribución de cola pesada, $P(X>x)\propto x^{-a}$. Aquí
$a=1.6$ para el tamaño de las empresas.
\item[Coeficiente de variación.] Desviación estándar dividida por la media. 8~\% y
25~\% son los coeficientes de variación de la frecuencia de pago.
\item[PD verdadera.] $\mathrm{logit}(\mathrm{PD}_i)=b_\text{tipo}+\sum_kc_kz_k(i)+
\sigma_uu_i$.
\item[$\mathrm{logit}(p)$.] $\ln\frac{p}{1-p}$, el logaritmo de la razón de
probabilidades.
\item[$z_k(i)$.] Variables clásicas estandarizadas (media 0, desviación 1) dentro de cada
tipo de cliente.
\item[$c_k$.] Pesos: endeudamiento $+0.7$, ingreso $-0.5$, antigüedad $-0.3$, edad
$-0.2$, formalidad $-0.5$.
\item[$u_i$, $\sigma_u$.] Componente latente $u_i\sim N(0,1)$ y su escala.
\item[Mora observada.] Número de moras en 12 meses $\sim\mathrm{Poisson}(m\,e^{au_i})$:
quien tiene $u_i$ alto tiende a tener más moras.
\item[$b_\text{tipo}$.] Intercepto calibrado por bisección para que la PD media sea 5~\%
en personas y 3~\% en empresas.
\item[Bisección.] Método numérico que encuentra la raíz de una función partiendo un
intervalo por la mitad repetidamente.
\item[SMMLV.] Salario mínimo mensual legal vigente. Unidad de los ingresos y montos.
\end{description}
\end{terminos}

\diapo{15}{Modelo dinámico multicapa}{1.5 min}
\begin{decir}
Esta es la versión del modelo para la cartera. Tiene la misma estructura que la ecuación
general con tres cambios.

Primero, el tiempo se mide en meses y el déficit como fracción del ingreso total del
cliente. Un déficit de 0.5 significa medio mes de ingreso. El ingreso total incluye lo que
llega de fuera de la red del banco.

Segundo, el contagio y el acople de fase tienen un valor por capa. La nómina transmite
más que las transferencias. Además, cada cliente ajusta su ritmo según la participación
de cada pagador en su ingreso: el empleador pesa mucho más que un amigo que le transfiere
dinero de vez en cuando.

Tercero, hay ruido aleatorio que representa choques individuales de liquidez, y hay
impagos idiosincráticos: cada cliente puede caer por razones propias, con una tasa que
corresponde a su probabilidad de impago.

El colchón depende del riesgo: un cliente con PD alta tiene un colchón más pequeño.

Con una sola capa y sin ruido, este modelo coincide con el de las tareas 1 a 4. Lo
verificamos numéricamente.
\end{decir}
\begin{terminos}
\begin{description}
\item[$\ell$.] Índice de capa: N, S o T.
\item[$W^\ell_{ij}$.] Flujo mensual de $j$ a $i$ en la capa $\ell$ [SMMLV/mes].
\item[$s_i$.] Ingreso total de $i$, incluido el externo a la red [SMMLV/mes].
\item[$k_i$.] Ingreso que llega por la red, $\sum_\ell\sum_jW^\ell_{ij}$.
\item[$\beta_\ell$.] Contagio por capa [1/mes]: N 0.5, S 0.4, T 0.3.
\item[$K_\ell$.] Acople de fase por capa [rad/mes]: N 0.8, S 1.0, T 0.3.
\item[$W^\ell_{ij}/k_i$.] Participación del pagador $j$ en el ingreso de red de $i$
(normalización por fila).
\item[$\gamma\frac{k_i}{s_i}(1-r_i)$.] Déficit por desfase, proporcional a la fracción del
ingreso que llega por la red. $\gamma=0.5$/mes.
\item[$\sigma\xi_i(t)$.] Ruido blanco de intensidad $\sigma=0.3$. Se integra con el
método de Euler–Maruyama.
\item[Euler–Maruyama.] Integrador de ecuaciones estocásticas:
$x_{n+1}=x_n+f(x_n)\Delta t+\sigma\sqrt{\Delta t}\,N(0,1)$.
\item[Colchón $x_i^{(c)}$.] $c\,(\mathrm{PD}_\text{ref}/\mathrm{PD}_i)^\kappa\varepsilon_i$
con $c=1.1$, $\mathrm{PD}_\text{ref}=0.04$, $\kappa=0.35$ y $\varepsilon_i$ lognormal.
\item[Impago idiosincrático.] Instante $\tau_i\sim\mathrm{Exp}(\lambda_i)$ con
$\lambda_i=-\ln(1-\mathrm{PD}_i)/12$ por mes, de modo que $P(\tau_i<12)=\mathrm{PD}_i$.
\item[Pérdida $\lambda=1$.] Perder un pagador que aporta todo el ingreso suma un mes de
ingreso al déficit.
\item[Matrices dispersas.] Almacenan solo las entradas distintas de cero; permiten
carteras grandes.
\end{description}
\end{terminos}
\begin{preguntas}
\pr{¿De dónde salen los valores de los parámetros?}{Son ilustrativos, no calibrados. Se
eligieron para que la tasa de impago y $R$ estén en rangos plausibles (impago cercano al
10~\%, $R$ entre 0.09 y 0.17). La diapositiva 22 muestra cómo cambia el resultado con
$\gamma$.}
\end{preguntas}

\diapo{16}{Del modelo a lo que observa el banco}{1 min}
\begin{decir}
Ahora el paso clave: pasar del modelo a lo que el banco realmente ve.

Cada vez que la fase de un pagador completa un ciclo, el pagador ejecuta sus pagos: a
sus empleados, a sus proveedores, y su cuota al banco. Agregamos ruido en la fecha y en el
monto. Quien cae en impago deja de pagar. El resultado es una tabla con fecha, pagador,
receptor, capa y monto: la misma estructura de un registro bancario.

Dividimos el tiempo en dos ventanas. De 0 a 6 meses, observamos: solo esas transacciones
alimentan las variables. De 6 a 18 meses, medimos si el cliente cae en impago; esa es la
etiqueta. Quien cae durante la observación se excluye.

Con las transacciones estimamos tres cosas: la red $\hat W$, la fase de cada cliente
$\hat\theta$ y su coherencia con los pagadores $\hat r$. La fase se estima por eventos:
entre dos pagos consecutivos, la fase avanza de manera lineal una vuelta completa.

La figura compara la fase real con la estimada para un empleado formal y un informal. La
estimación es buena: el error mediano en la red es menor al 5~\% y la correlación entre la
coherencia estimada y la real supera 0.8.
\end{decir}
\begin{terminos}
\begin{description}
\item[Registro de transacciones.] Tabla \texttt{(t, src, dst, layer, amount)}: instante,
pagador, receptor, capa y monto.
\item[Pago por ciclo.] El pagador $j$ paga a $i$ un monto $W_{ij}\cdot2\pi/\omega_j$ por
ciclo, de modo que el flujo mensual medio es $W_{ij}$.
\item[Ciclo observado.] Conjunto de pagos de una misma entidad separados por menos de
medio día; incluye la cuota al banco.
\item[Ruido de fecha.] $\pm1$ día, compartido por todos los pagos del mismo ciclo.
\item[Ventana de observación.] Meses $[0,6)$.
\item[Ventana de predicción.] Meses 6 a 18. Etiqueta $y=1$ si el cliente cae en ese
periodo.
\item[$\hat W_{ij}$.] Monto total pagado por $j$ a $i$ en la observación, dividido por
el periodo activo de $j$: $\min(t_\text{obs},\ \text{último ciclo}+\text{periodo medio})$.
Así, si $j$ dejó de pagar a mitad de la ventana, $\hat W_{ij}$ refleja lo que $i$ recibía
mientras $j$ pagaba, que es lo que perdió.
\item[Fase por eventos.] Pikovsky, Rosenblum y Kurths (2001):
$\hat\theta(t)=2\pi\bigl(k+\frac{t-t_k}{t_{k+1}-t_k}\bigr)$, donde $t_k$ es el instante del
$k$-ésimo ciclo de pagos.
\item[$\hat r_i(t)$.] $\frac1{\hat k_i}\sum_j\hat W_{ij}\cos(\hat\theta_j-\hat\theta_i)$,
calculada en una malla de 0.05 meses.
\item[Fuga de información.] Uso, al entrenar, de información que no estaría disponible al
predecir. El código rechaza transacciones posteriores a la observación.
\end{description}
\end{terminos}
\begin{preguntas}
\pr{¿Cómo se estima la fase de alguien que recibe pero casi no paga?}{Su fase se estima
con sus propios ciclos de pago, incluida la cuota al banco. Si no tiene ingreso por la
red, no hay desincronización medible y se le asigna $\hat r=1$.}
\end{preguntas}

\diapo{17}{Variables en tres niveles y protocolo}{1 min}
\begin{decir}
Con lo observado construimos variables en tres niveles, cada uno contiene al anterior.

El nivel 1 es el score clásico: ingreso, endeudamiento, edad, antigüedad, productos,
moras y formalidad.

El nivel 2 agrega la red estática: cuánto del ingreso llega por la red y por cuál capa,
qué tan concentrado está en pocos pagadores, la centralidad del cliente, y la salud de su
empleador y de sus pagadores. Esto corresponde a lo que ya explora la literatura.

El nivel 3 agrega la dinámica, que es el aporte del proyecto: la coherencia media y su
variabilidad, el tiempo que el cliente pasa desincronizado, el desfase con sus pagadores,
la diferencia de ritmo, la irregularidad de los ciclos y el ingreso de pagadores que
dejaron de pagar.

El protocolo tiene dos regímenes. Con red, los acoples están encendidos. En el régimen
nulo, apagamos el contagio, el desfase y la pérdida, pero mantenemos los mismos clientes,
la misma red y los mismos números aleatorios. Generamos 20 carteras por régimen y
validamos con GroupKFold por cartera. Comparamos los niveles por pares dentro de cada
partición y reportamos intervalos de confianza del 95~\%.
\end{decir}
\begin{terminos}
\begin{description}
\item[Concentración (HHI).] Índice de Herfindahl–Hirschman del ingreso,
$\sum_j(\hat W_{ij}/\hat k_i)^2$: vale 1 si todo viene de un pagador.
\item[PageRank.] Centralidad que asigna más importancia a quien recibe de entidades
importantes; factor de amortiguación 0.85. Se reescala para que la media sea 1.
\item[Salud del empleador y de los pagadores.] Tamaño, endeudamiento y moras del
empleador; moras y endeudamiento medios de los pagadores, ponderados por lo que pagan.
\item[$\hat r$ medio y desviación.] Promedio y desviación estándar de $\hat r_i(t)$ en la
observación.
\item[Tiempo desincronizado.] Fracción del tiempo con $\hat r_i<0.5$.
\item[Desfase.] Media en el tiempo de $|\arg\sum_j\hat W_{ij}e^{i(\hat\theta_j-
\hat\theta_i)}|$: el ángulo del desfase promedio con los pagadores, ponderado por lo que
paga cada uno, entre $0$ y $\pi$.
\item[Desajuste de ritmo.] $|\hat\omega_i-\bar\omega_\text{pagadores}|/2\pi$, en ciclos
por mes.
\item[Irregularidad de ciclos.] Coeficiente de variación de los intervalos entre ciclos,
propio y de los pagadores.
\item[Ingreso perdido.] Fracción del ingreso de red que venía de pagadores que dejaron de
pagar durante la observación.
\item[Régimen nulo.] $\beta_\ell=\gamma=\lambda=0$. Las fases y los pagos existen, pero no
crean déficit.
\item[Mismos sorteos.] La semilla de cada cartera genera cuatro flujos aleatorios
independientes (umbrales, impagos idiosincráticos, ruido y transacciones) con
\texttt{SeedSequence}. Los dos regímenes usan exactamente los mismos.
\item[Diferencia pareada.] Diferencia de métricas entre dos niveles evaluados en la misma
partición. Elimina la variación debida a qué carteras quedaron en cada partición.
\item[Intervalo de confianza 95~\%.] $\bar d\pm t_{0.975,4}\,s_d/\sqrt5$, con $\bar d$ la
media y $s_d$ la desviación de las 5 diferencias, y $t_{0.975,4}=2.78$ el cuantil de la
distribución $t$ de Student con 4 grados de libertad.
\end{description}
\end{terminos}
\begin{preguntas}
\pr{¿Por qué comparar contra la red estática y no solo contra el score?}{Porque superar
al score con información de red ya está documentado. La pregunta es si la dinámica añade
algo que la estructura no tiene.}
\end{preguntas}

% ===================================================================
\section{Resultados}

\diapo{18}{Validación numérica}{45 s}
\begin{decir}
Antes de interpretar resultados hay que confiar en el código. Cada componente se
comparó con un resultado conocido.

La integración coincide con la solución exacta del sistema lineal con error menor que
$10^{-6}$. La redistribución conserva el déficit total con precisión de $10^{-10}$. La
sincronización reproduce el valor teórico de Kuramoto. El modelo multicapa con una sola
capa coincide exactamente con la versión anterior. Y la coherencia estimada desde
transacciones se correlaciona por encima de 0.8 con la real.

El conjunto de datos se regenera idéntico bit a bit. En total el proyecto tiene 127
pruebas automáticas, y todas pasan.
\end{decir}
\begin{terminos}
\begin{description}
\item[RK4.] Método de Runge–Kutta de cuarto orden: evalúa la derivada cuatro veces por
paso y combina los resultados. Error global proporcional a $\Delta t^4$. Integrador de
las tareas 1–5.
\item[$e^{Mt}x_0$.] Solución exacta de $\dot x=Mx$ con condición inicial $x_0$; $e^{Mt}$
es la exponencial de una matriz.
\item[Varianza de Ornstein–Uhlenbeck.] Con solo $-\alpha x+\sigma\xi$, la varianza
estacionaria es $\sigma^2/(2\alpha)$. La simulación la reproduce dentro del 10~\% y
valida el ruido.
\item[Bit a bit.] Dos ejecuciones producen archivos idénticos byte por byte.
\item[Prueba automática.] Función que ejecuta una parte del código y comprueba su
resultado; se corre con \texttt{pytest}.
\end{description}
\end{terminos}

\diapo{19}{Carteras simuladas}{45 s}
\begin{decir}
Estas son las carteras del experimento. En el régimen nulo hay unos 40\,000 clientes
modelados, con una tasa de impago a doce meses de 5.9~\%. En el régimen con red hay unos
37\,000, porque más clientes caen durante la observación, y la tasa de impago sube a
9.6~\%. La diferencia es el efecto de la red. El índice de amplificación está entre 0.09
y 0.17: la red amortigua, no está cerca de la inestabilidad.

La figura muestra la coherencia estimada por grupo. Los empleados formales están
sincronizados con su nómina. Los informales reciben ingresos de pagadores con ritmos
distintos y su coherencia está cerca de cero. En el régimen con red, los impagos por
umbral, 2\,335, superan a los idiosincráticos, 1\,264.
\end{decir}
\begin{terminos}
\begin{description}
\item[Clientes modelados.] Clientes vivos al final de la observación, sumados en las 20
carteras.
\item[Diagrama de caja.] La caja va del primer al tercer cuartil; la línea interior es la
mediana; los bigotes muestran la dispersión.
\item[Impago por umbral.] Causado por la dinámica: el déficit superó el colchón.
\item[Impago idiosincrático.] Programado por la PD propia, independiente de la red.
\end{description}
\end{terminos}

\diapo{20}{Desempeño predictivo por nivel}{1 min}
\begin{decir}
Este es el resultado principal. En el régimen con red, con regresión logística, el AUC
pasa de 0.769 con el score clásico a 0.784 al agregar la red estática, y a 0.795 al
agregar la dinámica. Con XGBoost el patrón es el mismo y llega también a 0.795. Las
demás métricas mejoran en la misma dirección: la PR-AUC, que es la métrica relevante
cuando los impagos son pocos, sube de 0.277 a 0.323.

En el régimen nulo, en la parte izquierda de la figura, el AUC es más alto, 0.81, pero
los tres niveles son iguales. Es más alto porque allí casi todos los impagos dependen
solo de la probabilidad propia de cada cliente, que es justamente lo que mide el score
clásico. Los tres niveles son iguales porque la red no tiene nada que aportar.
\end{decir}
\begin{terminos}
\begin{description}
\item[Niveles 1, 2 y 3.] Clásico, clásico más red estática, y clásico más red estática
más red dinámica.
\item[PR-AUC.] Promedio de la precisión sobre todos los niveles de exhaustividad. Un
modelo al azar obtiene la tasa de impago, aquí cerca de 0.096.
\item[Precisión.] De los clientes marcados como riesgosos, fracción que cae.
\item[Exhaustividad (\emph{recall}).] De los clientes que caen, fracción marcada como
riesgosa.
\item[KS.] Estadístico de Kolmogorov–Smirnov: máxima distancia entre las distribuciones
acumuladas del score de quienes caen y de quienes no. Es una métrica habitual en banca.
\item[Brier.] Error cuadrático medio de la probabilidad, $\frac1n\sum(p_i-y_i)^2$.
Menor es mejor.
\item[ECE.] Error de calibración esperado: diferencia media entre la probabilidad
predicha y la frecuencia observada, por tramos de probabilidad.
\item[Fuera de muestra.] Métricas calculadas sobre carteras que el modelo no vio.
\end{description}
\end{terminos}
\begin{preguntas}
\pr{¿Por qué el AUC del régimen nulo es mayor?}{En el nulo el impago depende casi solo
de la PD de cada cliente, que es función de las variables clásicas. En el régimen con
red, parte del impago viene de la red y el score clásico no lo ve. Por eso el score
clásico pierde AUC con red y las variables de red recuperan una parte.}
\end{preguntas}

\diapo{21}{Mejoras pareadas y control nulo}{1 min}
\begin{decir}
Para saber si las diferencias son reales, comparamos los niveles dentro de cada
partición. La figura muestra cada mejora con su intervalo de confianza del 95~\%; en
verde, las que están por encima de cero.

En el régimen con red, la dinámica mejora a la red estática en 0.011 de AUC con
regresión logística, con intervalo entre 0.006 y 0.016, y en 0.014 con XGBoost. La
mejora es positiva en las cinco particiones. El intervalo no incluye el cero.

En el régimen nulo, la mejora está entre $-0.002$ y $+0.001$. Este es el control: cuando
el mecanismo no existe, el método no encuentra nada. Eso descarta que la mejora venga de
una fuga de información o de que el modelo tenga más variables.
\end{decir}
\begin{terminos}
\begin{description}
\item[$\Delta\mathrm{AUC}$.] AUC del nivel superior menos AUC del nivel de comparación,
en la misma partición.
\item[Gráfico de bosque.] Cada fila es una comparación: el punto es la media y la barra
el intervalo de confianza.
\end{description}
\end{terminos}
\begin{preguntas}
\pr{¿Una mejora de 0.01 en AUC es relevante?}{En riesgo de crédito sí. El Gini, que usan
los bancos, es $2\,\mathrm{AUC}-1$, así que la mejora es de 0.022 puntos de Gini sobre un
modelo que ya incluye la red estática. Sobre el score clásico, la mejora total con
XGBoost es de 0.031 de AUC. Además, el valor depende de parámetros ilustrativos: lo que
el experimento establece es que la mejora existe, es estable y se anula sin el
mecanismo.}
\pr{¿Más variables no siempre mejoran el AUC?}{No fuera de muestra: variables sin
información agregan ruido. El régimen nulo lo muestra: con las mismas variables, la
mejora es cero.}
\end{preguntas}

\diapo{22}{Respuesta a la intensidad del mecanismo}{1 min}
\begin{decir}
Este experimento adicional responde de dónde viene la mejora. Variamos $\gamma$, el
parámetro que convierte el desfase en déficit, y dejamos todo lo demás igual.

Con $\gamma$ igual a cero, el contagio y la pérdida de ingreso siguen activos, y aun así la
dinámica no aporta nada. La mejora aparece solo cuando el desfase genera déficit, y crece
de manera monótona: 0.001, 0.008 y 0.029 de AUC con $\gamma$ igual a 0.25, 0.5 y 1.

Conclusión: la información que aportan las variables dinámicas proviene de la
desincronización, no del contagio. Es una relación dosis-respuesta: más mecanismo, más
señal.
\end{decir}
\begin{terminos}
\begin{description}
\item[Sensibilidad.] Variación del resultado ante cambios en un parámetro.
\item[Diseño del experimento.] Regresión logística, 10 carteras por valor de $\gamma$,
misma validación; barras de error con IC 95~\%.
\item[Monótona.] Que siempre crece (o siempre decrece) al aumentar el parámetro.
\item[Dosis-respuesta.] Término de epidemiología: el efecto crece con la intensidad de la
causa. Es un criterio clásico para atribuir un efecto a una causa.
\end{description}
\end{terminos}

\diapo{23}{Interpretación y relación sincronía--impago}{1 min}
\begin{decir}
A la izquierda, la importancia de cada variable en el modelo híbrido. Las clásicas,
en gris, siguen siendo las más importantes: endeudamiento, ingreso y moras. Las de red
estática, en azul, y las dinámicas, en rojo, aportan información adicional. Entre las
dinámicas destacan la variabilidad de la coherencia, su nivel medio y el desfase. El
modelo híbrido complementa al score; no lo reemplaza.

A la derecha, la tasa de impago según la coherencia estimada. En el régimen con red, el
impago pasa de 4.4~\% en los clientes bien sincronizados a entre 11 y 14~\% en los
desincronizados. En el régimen nulo también hay una diferencia, más pequeña, porque los
informales son más riesgosos y además están menos sincronizados. Por eso esta gráfica
por sí sola no prueba nada: la prueba es la comparación de modelos de las diapositivas
anteriores.
\end{decir}
\begin{terminos}
\begin{description}
\item[SHAP.] \emph{SHapley Additive exPlanations} (Lundberg y Lee, 2017). Reparte la
predicción de un cliente entre sus variables usando los valores de Shapley de la teoría
de juegos: la contribución de cada variable promediada sobre todos los órdenes en que se
pueden agregar las variables. La suma de las contribuciones más un valor base da la
predicción.
\item[Importancia media $|\mathrm{SHAP}|$.] Promedio sobre los clientes del valor
absoluto de la contribución de cada variable, en unidades de log-odds.
\item[Log-odds.] $\ln\frac{p}{1-p}$, la escala en que el modelo suma contribuciones.
\item[Variable de confusión.] Variable que influye a la vez en la causa aparente y en el
efecto. Aquí, la informalidad.
\end{description}
\end{terminos}
\begin{preguntas}
\pr{¿El modelo está calibrado?}{Sí. En la curva de calibración ambos modelos siguen la
diagonal y el híbrido separa mejor el decil de mayor riesgo: 35~\% de impago observado
frente a 30~\% del clásico. La figura está en la guía completa y en la presentación v1.}
\end{preguntas}

\diapo{24}{Escenario: quiebra de un empleador}{1 min}
\begin{decir}
El modelo también sirve para pruebas de estrés. Simulamos la quiebra del mayor empleador
de la cartera, con 74 empleados, en el primer mes, y promediamos cinco semillas.

En el régimen con red, el 12~\% de los empleados cae de inmediato, porque pierde todo su
salario de un momento a otro. A los seis meses el impago acumulado de esos empleados es
13~\%, frente a 7~\% sin quiebra. En el régimen nulo la quiebra no tiene efecto, como
corresponde.

La aplicación directa para un banco es medir cuánto riesgo concentra en los empleados de
una misma empresa.
\end{decir}
\begin{terminos}
\begin{description}
\item[Prueba de estrés.] Simulación de un escenario adverso extremo para medir pérdidas.
\item[Choque de quiebra.] Impago forzado de la firma en $t_0=1$ mes: deja de pagar
nómina, proveedores y cuota.
\item[Semilla.] Número que inicializa el generador aleatorio; misma semilla, mismo
resultado.
\end{description}
\end{terminos}

\diapo{25}{Alcance e interpretación}{45 s}
\begin{decir}
Para cerrar los resultados, separo lo que el diseño respalda de lo que no.

Respalda tres afirmaciones. Si el mecanismo existe, es detectable con las transacciones
de que ya dispone un banco. La dinámica aporta información que la red estática no tiene.
Y el aporte se anula sin el mecanismo y crece con su intensidad.

No respalda que el mecanismo exista en una cartera real: los datos son sintéticos y los
parámetros ilustrativos. Además, supusimos que el banco ve toda la red de sus clientes.
\end{decir}
\begin{preguntas}
\pr{¿Qué haría falta para usarlo en un banco?}{Aplicar el mismo pipeline a una cartera
real anonimizada, bajo la Ley 1581 de 2012 de protección de datos personales. El código
ya recibe un registro de transacciones con la misma estructura.}
\end{preguntas}

% ===================================================================
\section{Implementación}

\diapo{26}{Arquitectura y reproducibilidad}{30 s}
\begin{decir}
El software está organizado en cuatro capas. La interfaz es una aplicación web. La capa
de servicio prepara las tablas que la interfaz muestra. El núcleo contiene el modelo y los
algoritmos. La configuración y los datos están separados del código. Cada capa solo
depende de la de abajo.

Todos los parámetros están en archivos de configuración, las semillas son explícitas, las
tablas tienen esquemas validados, hay 127 pruebas, integración continua y una imagen de
Docker que ejecuta la aplicación igual en cualquier equipo.
\end{decir}
\begin{terminos}
\begin{description}
\item[Streamlit.] Biblioteca de Python que convierte un programa en una aplicación web
interactiva.
\item[YAML.] Formato de texto para configuración, legible por personas.
\item[Esquema validado.] Lista de columnas, tipos y rangos que toda tabla debe cumplir;
si no se cumple, el programa se detiene.
\item[Integración continua.] Ejecución automática de pruebas en cada cambio del código
(GitHub Actions).
\item[Docker.] Herramienta que empaqueta el programa con todas sus dependencias en una
imagen que se ejecuta igual en cualquier computador.
\end{description}
\end{terminos}

\diapo{27}{Flujo de datos de la tarea 6}{30 s}
\begin{decir}
Este es el recorrido completo de la tarea 6: generar la cartera, simular el modelo,
producir las transacciones, estimar la red y las fases, construir las variables, validar
y comparar los modelos. Un solo comando ejecuta todo el flujo en unos 20 segundos.
\end{decir}
\begin{terminos}
\begin{description}
\item[\texttt{uv}.] Gestor de entornos y dependencias de Python.
\item[\texttt{pulso credit}.] Comando del proyecto que ejecuta el experimento de la
tarea 6 y guarda resultados y modelos finales.
\end{description}
\end{terminos}

\diapo{28}{Aplicación web}{30 s}
\begin{decir}
Todo lo presentado se puede explorar en la aplicación. A la izquierda, el simulador de
choques de las tareas 1 a 4: se elige una entidad, se aplica un choque y se ve la
cascada. A la derecha, la comparación de los tres niveles de la tarea 6 y la
probabilidad de impago por cliente, con su explicación. El manual de la aplicación
describe cada página.
\end{decir}

% ===================================================================
\section{Conclusiones}

\diapo{29}{Conclusiones}{1 min}
\begin{decir}
Primero, un único modelo dinámico integra contagio, redistribución, sincronización y
cascadas, y cada componente se valida contra resultados analíticos.

Segundo, el índice $R$ y el espectro del sistema dicen cuándo la red amortigua y cuándo
amplifica un choque.

Tercero, el aprendizaje automático predice el colapso del sistema con AUC de 0.98 en
milisegundos.

Cuarto, y es el resultado central: en una cartera multicapa, la sincronización de los
ciclos de pago, estimada solo con transacciones, mejora la predicción del impago sobre la
red estática. No tiene efecto en el control nulo y crece con la intensidad del
mecanismo.

Quinto, el software es reproducible, está probado y funciona como aplicación web.
\end{decir}

\diapo{30}{Trabajo futuro}{30 s}
\begin{decir}
El paso más importante es la validación con datos reales de un banco, anonimizados.
También queda pendiente una revisión sistemática de la literatura. En el modelo, faltan
las fechas de cuota como ritmo impuesto desde fuera, la retroalimentación del déficit
sobre el ritmo de pago, y el caso en que un banco solo ve a sus clientes, frente al
operador de Bre-B, que ve la red completa.
\end{decir}
\begin{terminos}
\begin{description}
\item[Revisión sistemática.] Búsqueda bibliográfica con protocolo explícito en bases como
Scopus y Web of Science.
\item[Forzamiento externo.] Ritmo impuesto desde fuera del sistema, por ejemplo la fecha
fija de la cuota.
\item[Observabilidad parcial.] Situación en la que solo se ve una parte de la red.
\end{description}
\end{terminos}

\diapo{31}{Referencias principales}{15 s}
\begin{decir}
Estas son las referencias principales. La bibliografía completa está en el repositorio.
Muchas gracias; quedo atento a sus preguntas.
\end{decir}

% ===================================================================
\section{Preguntas generales para la defensa}

Estas preguntas no dependen de una diapositiva concreta.

\begin{preguntas}
\pr{¿Cuál es la contribución original del proyecto?}{Mostrar, con un experimento
controlado, que la sincronización de los ciclos de pago, estimada solo con
transacciones, contiene información sobre el impago que no está en el score clásico ni
en la estructura estática de la red. Incluye el método para estimarla (fase por eventos y
coherencia local) y un control nulo que descarta artefactos.}
\pr{¿Por qué un físico trabaja en riesgo de crédito?}{Porque las herramientas son de la
física: sistemas dinámicos, análisis espectral, osciladores acoplados, transiciones de
fase y avalanchas. La red de pagos es un sistema complejo y estas herramientas permiten
responder cuándo amplifica y cómo se sincroniza.}
\pr{¿El resultado es causal o solo predictivo?}{En el mundo simulado la causa se conoce
porque la pusimos: el desfase genera déficit. El experimento de $\gamma$ muestra que la
señal depende de esa causa. Con datos reales, el método sería predictivo; establecer
causalidad requeriría otro diseño.}
\pr{¿Cómo se evitó la fuga de información?}{De cuatro formas. Las variables solo usan
transacciones anteriores al fin de la observación, y el código rechaza cualquier pago
posterior. Las columnas de verdad del modelo están vetadas por esquema. La validación
separa carteras completas. Y el control nulo muestra mejora cero.}
\pr{¿Por qué no se usó una GCN en la tarea 6?}{Las variables de red ya resumen la
información relevante de los vecinos de cada cliente, y la pregunta es qué tipo de
información aporta. Comparar niveles de variables con el mismo modelo responde esa
pregunta de forma más limpia que un modelo de grafo de caja negra.}
\pr{¿Qué pasa si los parámetros son otros?}{El tamaño de la mejora cambia: crece con
$\gamma$ y se anula sin él. Lo que se mantiene es el patrón: sin mecanismo no hay
mejora, con mecanismo sí. Los parámetros están en \texttt{configs/credit.yaml} y el
experimento completo se vuelve a correr en unos 20 segundos.}
\pr{¿Funciona si el banco no ve toda la red?}{No se probó. Un banco ve solo a sus
clientes y los pagos que pasan por él. Es parte del trabajo futuro: el operador de un
sistema interoperable como Bre-B tendría la vista completa.}
\pr{¿Cuánto tarda todo?}{Las pruebas rápidas, unos 10 segundos. El experimento de la
tarea 6, unos 20 segundos. El conjunto de datos de la tarea 5, unos 20 minutos en ocho
núcleos.}
\end{preguntas}

% ===================================================================
\section{Glosario de siglas, modelos y métodos}

\begin{longtable}{p{0.2\textwidth}p{0.74\textwidth}}
\toprule
\textbf{Término} & \textbf{Significado} \\
\midrule
\endhead
AUC & Área bajo la curva ROC. Probabilidad de que un positivo al azar reciba más riesgo
que un negativo al azar. \\
BA & Barabási–Albert: red con conexión preferencial y nodos concentradores. \\
Bre-B & Sistema de pagos inmediatos del Banco de la República de Colombia. \\
Brier & Error cuadrático medio de una probabilidad predicha. \\
DANE & Departamento Administrativo Nacional de Estadística. \\
DebtRank & Medida de propagación de estrés en redes financieras (Battiston et al.,
2012). \\
ECE & Error de calibración esperado. \\
Eisenberg–Noe & Modelo de liquidación de obligaciones en una red (2001). \\
ER & Erdős–Rényi: red aleatoria con probabilidad de conexión uniforme. \\
Euler–Maruyama & Integrador de ecuaciones diferenciales con ruido. \\
Fiedler & Vector propio de la conectividad algebraica $\lambda_2$; separa comunidades. \\
GCN & Red neuronal convolucional sobre grafos. \\
Gini & $2\,\mathrm{AUC}-1$. Métrica de discriminación usada en banca. \\
GroupKFold & Validación cruzada que mantiene cada grupo en una sola partición. \\
HHI & Índice de Herfindahl–Hirschman: concentración, suma de participaciones al
cuadrado. \\
IC 95~\% & Intervalo de confianza del 95~\%. \\
Kuramoto & Modelo de osciladores de fase acoplados. \\
KS & Estadístico de Kolmogorov–Smirnov: máxima separación entre distribuciones
acumuladas. \\
Laplaciano & Matriz $D-W$ que describe flujos que se conservan en una red. \\
Lorentz & Distribución de Cauchy; da la solución exacta de Kuramoto en la red
completa. \\
NP & Núcleo-periferia. \\
Parquet & Formato de archivo para tablas, comprimido y con tipos fijos. \\
Pareto & Distribución de cola pesada: pocos valores muy grandes. \\
PD & Probabilidad de incumplimiento a un horizonte (aquí 12 meses). \\
Perron–Frobenius & Teorema sobre el valor propio dominante de matrices no negativas. \\
PINN & Red neuronal informada por la física. \\
PR-AUC & Área bajo la curva precisión-exhaustividad. \\
PageRank & Centralidad por recorridos aleatorios sobre la red. \\
$R$ & Índice de amplificación $\beta\rho(W)/\alpha$. \\
$r$, $r_i$ & Parámetro de orden global y coherencia local. \\
RK4 & Runge–Kutta de cuarto orden. \\
ROC & Curva de tasa de verdaderos positivos contra falsos positivos. \\
SBM & Modelo de bloques estocástico: la probabilidad de conexión depende del grupo de
cada nodo. Se usa en las transferencias. \\
SeedSequence & Generador de semillas independientes de NumPy. \\
SHAP & Explicación aditiva de predicciones basada en valores de Shapley. \\
SMMLV & Salario mínimo mensual legal vigente. \\
Spearman & Correlación entre rangos. \\
Watts & Modelo de cascadas por umbral fraccional (2002). \\
XGBoost & Árboles de decisión con \emph{gradient boosting}. \\
YAML & Formato de texto para archivos de configuración. \\
\bottomrule
\end{longtable}


\end{document}
